\chapter{Istniejące algorytmy planowania czasu procesora w nowoczesnych systemach}
\label{ch:wprowadzenie}

\subsection{System Linux i architektura EEVDF}
W nowoczesnych wersjach jądra systemu Linux, domyślnym algorytmem planowania (ang. \textit{scheduler}) dla procesów standardowych jest algorytm Earliest Eligible Virtual Deadline First (EEVDF). Zastąpił on wieloletni standard znany jako Completely Fair Scheduler (CFS), rozwiązując fundamentalny problem balansu między sprawiedliwością przydziału czasu a opóźnieniami w obsłudze zadań interaktywnych \cite{1}, \cite{38}. 

Działanie algorytmu EEVDF opiera się na ścisłych podstawach matematycznych \cite{15} i operuje na kilku kluczowych wartościach przypisywanych każdemu zadaniu:
\begin{itemize}
	\item \textbf{Żądany kwant czasu} (\textit{requested time slice / slice})
	\item \textbf{Wirtualny czas wykonania} (\textit{virtual runtime / vruntime})
	\item \textbf{Opóźnienie} (\textit{lag})
	\item \textbf{Priorytet} (\textit{priority / weight})
\end{itemize}

Kiedy nowe lub powracające ze stanu uśpienia zadanie trafia do kolejki procesów gotowych, otrzymuje bazowy kwant czasu (\textit{slice}). Wartość ta oparta jest na zmiennych konfiguracyjnych systemu (np. \texttt{sysctl\_sched\_base\_slice}) i jest dynamicznie skalowana na podstawie wagi procesu (wynikającej z jego priorytetu \textit{nice}) \cite{19}, \cite{26}. Następnie, gdy zadanie zaczyna być wykonywane przez procesor, jego wirtualny czas wykonania (\texttt{vruntime}) stale rośnie. Wartość ta odzwierciedla rzeczywisty czas spędzony na procesorze, znormalizowany przez wagę zadania. Zliczanie tego czasu opiera się na precyzyjnych licznikach sprzętowych procesora (np. w architekturze amd64), jednak sama aktualizacja wirtualnego czasu i struktur danych jest w całości realizowana programowo przez jądro systemu operacyjnego, zazwyczaj podczas obsługi przerwań zegarowych \cite{18}, \cite{19}.

Na bazie początkowo przypisanego kwantu czasu oraz \texttt{vruntime}, system wylicza wirtualny termin zakończenia (\textit{virtual deadline}). Reprezentuje on moment, do którego proces powinien w idealnych warunkach skonsumować swój przydzielony czas. Algorytm na bieżąco oblicza również opóźnienie (\textit{lag}), będące różnicą pomiędzy czasem, jaki zadanie powinno było otrzymać w idealnym, proporcjonalnym modelu podziału czasu (tzw. \textit{fluid-flow model}), a czasem, który faktycznie wykorzystało. Jeżeli wartość \textit{lag} jest mniejsza od zera, oznacza to, że proces otrzymał więcej czasu procesora niż pozostałe procesy rywalizujące o czas procesora. W takim scenariuszu proces traci kwalifikowalność (\textit{ineligible}) i jest kategorycznie wykluczony z możliwości przydziału sprzętowego do momentu, aż ogólny upływ czasu w systemie sprawi, że jego \textit{lag} ponownie wzrośnie do wartości nieujemnej \cite{1}, \cite{15}.

\subsection{Wybór procesu i struktury danych}
Działanie każdego planisty w nowoczesnych systemach operacyjnych opiera się na sprzętowym mechanizmie przerwań czasowych (ang. \textit{timer interrupts}). Ponieważ procesor nie może jednocześnie wykonywać własnego kodu schedulera oraz kodu aplikacji użytkownika, wbudowany kontroler przerwań z dużą częstotliwością wstrzymuje wykonywane zadanie, przekazując sterowanie do jądra systemu \cite{13}. Daje to algorytmowi możliwość przeliczenia wskaźników i, w razie potrzeby, wywłaszczenia procesu.

Scheduler EEVDF przechowuje informacje o aktywnych procesach w zrównoważonym drzewie czerwono-czarnym (ang. \textit{Augmented Red-Black Tree}). W przeciwieństwie do naiwnego podejścia, drzewo to nadal sortuje węzły na podstawie wirtualnego czasu wykonania (\texttt{vruntime}), co układa zadania od lewej do prawej. Kluczową innowacją jest jednak rozszerzenie (\textit{augmentation}) każdego węzła o dodatkowe pole \texttt{min\_deadline}, które przechowuje najkrótszy termin zakończenia z danego węzła i wszystkich jego potomków. Dzięki temu struktura działa jednocześnie jako drzewo wyszukiwania binarnego oraz jako kopiec (min-heap) dla wartości \textit{deadline}. Architektura ta pozwala systemowi w optymalnym czasie $O(\log n)$ odrzucić całe poddrzewa i natychmiast odnaleźć zadanie o najwcześniejszym \textit{virtual deadline}, które jednocześnie spełnia warunek kwalifikowalności ($lag \ge 0$) i to ono zostaje wybrane do działania na procesorze \cite{4}, \cite{19}.

\subsection{Mechanizm Run to Parity}
Algorytm EEVDF w swojej czystej, teoretycznej postaci zakłada wyidealizowany model maszyny obliczeniowej. W rzeczywistości fizyczne procesory opierają swoją wydajność na hierarchii pamięci podręcznej (Cache L1/L2/L3), która przechowuje dane powiązane z aktualnie wykonywanym procesem. Niezmodyfikowany mechanizm doboru zadań oparty wyłącznie na \textit{deadline'ach} prowadziłby do zjawiska tzw. \textit{over-scheduling} – zbyt częstych zmian aktualnie działającego zadania. Skutkowałoby to znacznymi stratami cykli procesora na ciągłą wymianę linii pamięci podręcznej (\textit{cache miss}) oraz narzutem związanym z przełączaniem kontekstu.

Aby rozwiązać ten problem, w systemie Linux wprowadzono mechanizm \textit{Run to Parity} \cite{26}. Zgodnie z tą regułą, proces, który już przejął kontrolę nad jednostką obliczeniową, z reguły nie zostaje wywłaszczony przez inne, standardowe zadania posiadające wcześniejszy \textit{virtual deadline}, dopóki jego własne opóźnienie (\textit{lag}) nie spadnie do zera. Dopiero gdy proces "wyrówna" swój dług wobec systemu (osiągając parytet), zdejmuje się z niego ochronę, pozwalając na przejęcie procesora przez kolejne kwalifikujące się zadanie. Mechanizm ten gwarantuje dłuższe, nieprzerwane odcinki pracy pojedynczego wątku, co maksymalizuje wykorzystanie pamięci podręcznej i zwiększa ogólną przepustowość sprzętu (throughput), nie naruszając jednocześnie długoterminowych matematycznych gwarancji sprawiedliwości algorytmu \cite{38}.